\phantomsection 
\addcontentsline{toc}{section}{LỜI MỞ ĐẦU}
\section*{LỜI MỞ ĐẦU}
Trong thời đại công nghệ 4.0 hiện nay, việc ứng dụng phân tích dữ liệu và trí tuệ nhân tạo vào các lĩnh vực đời sống, đặc biệt là quản lý tài nguyên thiên nhiên và bảo vệ sức khỏe cộng đồng, đã trở thành một xu hướng tất yếu. Nước sạch là yếu tố sống còn, tuy nhiên việc xác định chất lượng nước một cách nhanh chóng và chính xác vẫn là thách thức lớn đối với nhiều khu vực.

Ngành khoa học dữ liệu nói chung và kho dữ liệu nói riêng cung cấp các công cụ mạnh mẽ để lưu trữ, tích hợp và phân tích khối lượng lớn các chỉ số hóa lý của nước, từ đó giúp các chuyên gia đưa ra những quyết định kịp thời về mức độ an toàn của nguồn nước.

Xuất phát từ yêu cầu thực tiễn đó, đề tài \textbf{Phân tích và xây dựng mô hình dự đoán độ an toàn của nước (Water Potability Analysis)} từ tập dữ liệu \href{https://www.kaggle.com/datasets/uom190346a/water-quality-and-potability}{\color{blue} Kaggle dataset} tập trung ứng dụng các quy trình ETL hiện đại và các mô hình học máy để phân loại nguồn nước. Mã nguồn dự án được lưu trữ công khai \href{https://github.com/Shootingstar2608/data_warehouse_assignment}{\color{blue} tại đây.}

Mục tiêu của đề tài là xây dựng một kho dữ liệu hoàn chỉnh từ các nguồn dữ liệu thô ban đầu, thực hiện chuẩn hóa dữ liệu thông qua quy trình ETL phức tạp, và cuối cùng là huấn luyện các mô hình học máy để dự đoán khả năng uống được (potability) của nước dựa trên các thành phần như nồng độ pH, độ cứng, nồng độ Sulfate và các chỉ số hóa học khác.